\scanpage{9}\mspage{\handwritten{8}}

On doit trouver une approximation meilleure d'un cran à la classification topologique en prenant des $(2\text{-}\mathrm{Cat})$-groupes resp. des $(2\text{-}\mathrm{Cat})$-\unsure{groupes} abéliens, ou plus généralement des `` Gr-2-catégories '' resp. ``Grab-2-catégories'' . (NB U\deleted{\unsure{p}}ne Grab-catégorie défi\deleted{\unsure{b}}nit une Grab-\unsure{2catégorie} comme une catégorie définit une 2-catégorie; du point de vue semi-simplicial, cette opération serait analogue à celle de ``cosquélette\deleted{\unsure{x|r|\({}^{\times}\)}} tronqué''; Mais également une Grab-2-catégorie défini\deleted{e}\handwritten{t} \deleted{en}\handwritten{une} Grab-catégorie\deleted{\unsure{x|\({}^{\times}\)}} , savoir la catégorie des automorphismes de l'objet unité; cette \deleted{dé} opération serait analogue à celle de ``tuage des groupes d'homotopie en dim $\handwritten{>1}$''). Pour la notion gro\deleted{\unsure{t}}ssière ( Gr\deleted{ab} ou Grab) l'idée de ma lettre à Verdier était \unsure{qu'on} doit trouver la même chose que les tronqués \inserted{au cran 2}, dans la catégorie \uline{\unsure{homotopique}} des H-espaces - i.e. ce qu'on déduit des H-espaces par tuage des \unsure{groupes} d'homotopie en dimensions $\handwritten{>}2$.

\margin{\handwritten{\(\Vert\)}}% parallel strokes beside the question
(Demander à Quillen relations exactes entre \unsure{Objets|Hbjets|objets}-groupes et objets \deleted{ponctués} connexes de la catégorie homotopique \inserted{ponctuée}: forment-ils des catégories équivalentes par le foncteur espace des lacets\handwritten{?}).
% slanted
\margin{\handwritten{non, on a un foncteur \unsure{..|tr.|pl.} fidèle (?) mais \unsure{pas} \unsure{ess.t|ess.} surjectif...}}

Pour les \uline{vrais} groupes \inserted{abéliens} de $(2\text{-}\mathrm{Cat})$, l'approche de Quillen consiste à les identifier d'abord à des groupes abéliens de la catégorie $(\mathrm{Bicat})$, donc à des grpupes abéliens bi-semi\deleted{\unsure{e}}simpliciaux particuliers (conditions sur chaque ligne et chaque colonne). Appliquant la version bi-semisimpliciale du tapis de Dold, on trouve que \deleted{cette ca} la catégorie de ces objets est équivalente à la catégorie des \deleted{objets} bicomplexes \inserted{de chaines} de groupes \unsure{abéliens} qui sont nuls en dehors du carré $\deleted{\unsure{D}}0\handwritten{\leq}i,j\handwritten{\leq}1$. \deleted{Quillen dit que} La \deleted{colonne}\inserted{\handwritten{cond.}} que la première ligne corresponde à un \deleted{ensemble} groupe semisimplicial trivial signifie alors qu'elle n'a que le terme de degré zéro, i.e. que le bicomplexe est réduit à

% Arrow strokes in the bicomplex are in pen.
\diagram{\begin{tikzcd}[ampersand replacement=\&,row sep=small,column sep=small]
 C_{oo} \& 0\arrow[l] \& \handwritten{0}\arrow[l]\\
 \deleted{\unsure{C|0}}C_{1o}\arrow[u] \& C_{11}\arrow[l]\arrow[u] \& \handwritten{0}\arrow[l]\\
 \handwritten{\unsure{0}}\arrow[u] \& \&
\end{tikzcd}}

% archive sheet number: 8
